% =====================================================================
% Required preamble (place the following in your main .tex file):
%
%   \usepackage{tikz}
%   \usepackage{graphicx}   % for \resizebox
%   \usepackage{xcolor}     % for \definecolor (loaded by tikz, listed for clarity)
%   \usetikzlibrary{positioning, fit, arrows.meta, backgrounds}
%
%   \definecolor{somblue}{RGB}{207,226,243}    % intra-case  (soft blue)
%   \definecolor{somgreen}{RGB}{212,237,218}   % resource    (soft green)
%   \definecolor{sompeach}{RGB}{252,224,199}   % inter-case  (soft peach)
% =====================================================================

\begin{figure}[ht]
\centering
\resizebox{\textwidth}{!}{%
\begin{tikzpicture}[
    >={Latex[length=2.2mm]},
    node distance=0.9cm,
    every node/.style={font=\sffamily},
    box/.style={
        draw, line width=0.4pt, rounded corners=2pt,
        align=center, text width=4.2cm,
        inner sep=4pt, minimum height=0.85cm
    },
    input/.style={box, fill=gray!12, text width=5cm},
    intra/.style={box, fill=somblue},
    res/.style  ={box, fill=somgreen},
    inter/.style={box, fill=sompeach},
    merge/.style={box, fill=gray!12, text width=10cm},
    arrow/.style={->, line width=0.6pt},
    grouplabel/.style={
        font=\sffamily\bfseries, align=center,
        text width=4.4cm, inner sep=4pt
    },
    grouprect/.style={
        draw=gray!70, dashed, rounded corners=5pt, inner sep=8pt
    },
]

% ----- Input -----
\node[input] (input) at (0,0) {XES Event Log $L$};

% ----- Branch headers (column anchors) -----
\node[grouplabel] (h_intra) at (-5.8,-1.7)
    {Intra-case branch\\{\normalfont\small\itshape unit: events within a case}};
\node[grouplabel] (h_res)   at ( 0   ,-1.7)
    {Resource branch\\{\normalfont\small\itshape unit: resource $\times$ time bin $\Delta t$}};
\node[grouplabel] (h_inter) at ( 5.8,-1.7)
    {Inter-case branch\\{\normalfont\small\itshape unit: global time bin $\Delta t$}};

% ----- Branch 1: Intra-case -----
\node[intra, below=of h_intra] (i1) {Per-case event windows\\(size $w^{*}$)};
\node[intra, below=of i1]      (i2) {Feature vectors: activity one-hot,\\time deltas, position};
\node[intra, below=of i2]      (i3) {PCA $\to d^{*}_{\text{intra}}$};
\node[intra, below=of i3]      (i4) {SOM$_{\text{intra}}$ (e.g.\ $4\times 4$)};
\node[intra, below=of i4]      (i5) {Per-case state sequence};

% ----- Branch 2: Resource -----
\node[res, below=of h_res] (r1) {Resource-bin slicing (width $\Delta t$)};
\node[res, below=of r1]    (r2) {Feature vectors: activity histogram,\\workload, handover, waiting time};
\node[res, below=of r2]    (r3) {PCA $\to d^{*}_{\text{res}}$};
\node[res, below=of r3]    (r4) {SOM$_{\text{res}}$ (e.g.\ $3\times 3$)};
\node[res, below=of r4]    (r5) {Per-resource state sequence};

% ----- Branch 3: Inter-case -----
\node[inter, below=of h_inter] (n1) {Global time-bin slicing (width $\Delta t$)};
\node[inter, below=of n1]      (n2) {Feature vectors: activity \& DFG-edge\\histograms, active cases, gaps};
\node[inter, below=of n2]      (n3) {PCA $\to d^{*}_{\text{inter}}$};
\node[inter, below=of n3]      (n4) {SOM$_{\text{inter}}$ (e.g.\ $3\times 3$)};
\node[inter, below=of n4]      (n5) {Global state sequence};

% ----- Dashed grouping rectangles (drawn behind the branches) -----
\begin{scope}[on background layer]
  \node[grouprect, fit=(h_intra)(i1)(i5)] {};
  \node[grouprect, fit=(h_res)  (r1)(r5)] {};
  \node[grouprect, fit=(h_inter)(n1)(n5)] {};
\end{scope}

% ----- Merged bottom section (centered under resource column) -----
\node[merge, below=1.7cm of r5] (m1) {Per-SOM frequency time series \ $+$\ joint state signal};
\node[merge, below=of m1]       (m2) {PELT change-point detection};
\node[merge, below=of m2]       (m3) {Detected drift points};
\node[merge, below=of m3]       (m4) {Evaluation vs.\ CDLG ground truth\\(precision, recall, F1, mean delay)};

% ----- Arrows: input -> first step of each branch -----
\draw[arrow] (input.south) -- (i1.north);
\draw[arrow] (input.south) -- (r1.north);
\draw[arrow] (input.south) -- (n1.north);

% ----- Arrows: within each branch -----
\foreach \a/\b in {%
    i1/i2, i2/i3, i3/i4, i4/i5,
    r1/r2, r2/r3, r3/r4, r4/r5,
    n1/n2, n2/n3, n3/n4, n4/n5}
  \draw[arrow] (\a) -- (\b);

% ----- Fan-in: state sequences -> merge node -----
\draw[arrow] (i5.south) -- (m1.north);
\draw[arrow] (r5.south) -- (m1.north);
\draw[arrow] (n5.south) -- (m1.north);

% ----- Arrows: merge chain -----
\draw[arrow] (m1) -- (m2);
\draw[arrow] (m2) -- (m3);
\draw[arrow] (m3) -- (m4);

\end{tikzpicture}%
}
\caption{Three-branch pipeline for state-based concept drift detection in
traditional event logs. Each branch independently encodes one perspective
(intra-case, resource, inter-case), compresses with PCA, and discovers
states via a Self-Organizing Map. State frequencies form time series whose
change points (PELT) yield drift candidates.}
\label{fig:pipeline}
\end{figure}
